\chapter{Теория}

Линейная и параболическая модели являются частными случаями более
сложной модели – полиномиальной. Построить модель регрессии – это значит
найти параметры той функции, которая будет в ней фигурировать. Для
линейной регрессии – два параметра: коэффициент и свободный член.
Полиномиальная регрессия может применяться в математической
статистике при моделировании трендовых составляющих временных рядов.
Временной ряд – это, по сути, ряд чисел, которые зависят от времени.
Например, средние значения температуры воздуха по дням за прошедший год,
или доход предприятия по месяцам. Порядок моделируемого полинома
оценивается специальными методами, например, критерием серий. Цель
построения модели полиномиальной регрессии в области временных рядов всё
та же – прогнозирование.

\endinput